\appendix
\chapter[Apêndice - IMAGENS E MODELOS DE CORES]{Imagens e modelos de cores} \label{sec:modeloscores}
%\setlength{\afterchapskip}{-\baselineskip}
\pagestyle{fancy}



Existem diversos modelos em uso atualmente, os principais são o \siglai{RGB}{\textit{Red, Green, Blue} – vermelho, verde, azul}, o \siglai{CMYK}{\textit{Cyan, Magenta, Yellow, Black} – ciano, magenta, amarelo, preto} e o \siglai{HSI}{\textit{Hue, Saturation, Intensity} – matiz, saturação, intensidade}. A Figura \ref{fig:modeloscores1} mostra o modelo \sig{RGB} e o modelo \sig{CMYK}, o primeiro é baseado em cores primárias em um formato aditivo, conforme a Figura \ref{fig:modeloscores1}(c), uma imagem colorida \sig{RGB} é obtida a partir da sobreposição das matrizes de intensidade vermelha, verde e azul. Já o modelo \sig{CMYK} possui formato subtrativo e usa cores secundárias,  neste modelo uma cor absorve uma cor primária e reflete as outras duas. 
\afterpage{
\clearpage
\def\firstcircle{ (90:0.55) circle (1cm)}
\def\secondcircle{ (210:0.55) circle (1cm)}
\def\thirdcircle{  (330:0.55) circle (1cm)}
\def\firstcirclergb{ ( 90:0.55) circle (1)}
\def\secondcirclergb{(210:0.55) circle (1)}
\def\thirdcirclergb{(330:0.55) circle (1)}
\begin{figure}
\centering
\caption{Modelos de cores.} \label{fig:modeloscores1}
\begin{subfigure}[b]{.3 \linewidth}
\centering
\begin{tikzpicture}[blend group=screen]
\begin{scope}
\fill[red] \firstcircle;
\fill[green] \secondcircle;
\fill[blue] \thirdcircle;
\end{scope}
%\fill[red!90!white]
  \begin{scope}
  \clip \firstcirclergb;
  \clip \secondcirclergb;
  \fill [white] \thirdcirclergb;
  \end{scope}
\end{tikzpicture}
\subcaption{ \label{fig:rgb}}
\end{subfigure}
\qquad 
\begin{subfigure}[b]{0.3\linewidth}
\centering
\begin{tikzpicture}
\begin{scope}
\fill[yellow] \secondcircle;
\fill[magenta] \firstcircle;
\fill[cyan] \thirdcircle;
\end{scope}

\begin{scope}
\clip \firstcircle;
\fill[red]
\secondcircle;
\end{scope}

\begin{scope}
\clip \secondcircle;
\fill [green] \thirdcircle;
\end{scope}

\begin{scope}
\clip \firstcircle;
\fill [blue] \thirdcircle;
\end{scope}

\begin{scope}
\clip \firstcircle;
\clip \secondcircle;
\fill [black] \thirdcircle;
\end{scope}
\end{tikzpicture}
\subcaption{ \label{fig:cmyk}}
\end{subfigure}
\qquad

\begin{subfigure}[b]{0.2\linewidth}
\centering
\input{figuras/cap2/planosrgb} 
\subcaption{\label{fig:rgbsobreposicao}}
\end{subfigure}
\caption*{\subref{fig:rgb} cores \sig{RGB}. \subref{fig:cmyk} cores \sig{CMYK}. \subref{fig:rgbsobreposicao} Sobreposição das matrizes de intensidade para gerar modelo de cores \sig{RGB}.}
\end{figure}}
\afterpage{
\begin{figure}[b]
\centering
\caption{Modelos de cores e dispositivos}\label{fig:modeloscores}
\includegraphics[scale=0.84]{figuras/cap2/modeloscores.jpg} 
\caption*{\textbf{(a)} Subespaço de cores correspondente ao espectro visível determinado pelo \sig{CIE} 1976. \textbf{(b)} Subespaços \sig{RGB} e \sig{CMYK} em relação ao espectro visível. \textbf{(c)} Relação de dispositivos de entrada e saída com os modelos de cores. \textbf{Fonte:} blog corisectelmo - Disponível em \url{http://corisectelmo.blogspot.com.br/2011/01/aula-21-espaco-de-cor-lab.html} acessado em ago. 2015.}
\end{figure}

%\fontewwwfoot{http://corisectelmo.blogspot.com.br/2011/01/aula-21-espaco-de-cor-lab.html}{ago. 2015}{\value{footnote}}
\clearpage
}

